\section{Method}

\subsection{Problem Formulation}

Let \(\mathcal{D}_s\) and \(\mathcal{D}_t\) denote a source domain and a target
domain. Each domain contains user--item interaction records
\((u,i,r,e)\), where \(u\) is a user, \(i\) is an item, \(r\) is an observed
rating, and \(e\) is a natural-language explanation or review-derived
explanation target. Given source-domain evidence and target-domain evidence,
the goal is to predict a target-domain rating \(\hat r_{ui}\) and generate a
target-domain explanation \(\hat e_{ui}\) for a user--item pair.

The target explanation should satisfy two constraints. First, it should be
content-grounded: the explanation should describe evidence relevant to the
user and item rather than generic domain text. Second, it should be
domain-appropriate: the explanation should use target-domain expression style
without allowing style to replace recommendation evidence. ODCR formalizes
these constraints through content evidence \(C\), dynamic style \(S^{(t)}\),
and counterfactual reliability \(R_{cf}\).

\subsection{Dynamic Feedback View of Explanation Degeneration}

D4C diagnoses explanation degeneration as a causal shortcut involving hidden
textual attributes \(W\) \cite{d4c}. In the static reading, \(W\) captures
domain-level textual factors that can affect users, items, and explanations.
The higher-value reading for ODCR is dynamic: textual attributes and generated
explanations can reinforce one another over time,
\[
  W^{(t)} \rightarrow E^{(t)} \rightarrow W^{(t+1)} .
\]
When this loop is left uncontrolled, a generator can learn the shortcut
\(W^{(t)} \rightarrow E^{(t)}\) instead of the intended interaction path
\((U,I) \rightarrow E^{(t)}\). The resulting explanation may fit the target
domain's surface language while failing to explain the specific user--item
preference.

This feedback view motivates a control problem. The model should preserve the
part of textual evidence that supports the recommendation, allow style to adapt
across domains, and prevent unreliable counterfactual shifts from being treated
as trustworthy training evidence. ODCR keeps D4C's causal diagnosis but
replaces the coarse textual shortcut with explicit control variables.

\subsection{From Coarse Textual Attributes to ODCR Control Variables}

ODCR refines the D4C textual attribute at time \(t\) as
\[
  W^{(t)} \Rightarrow (C, S^{(t)}, R_{cf}) .
\]
Here \(C\) is stable content evidence: it captures what should be explained and
what can support rating prediction. \(S^{(t)}\) is a dynamic style state: it
captures how content is expressed in a domain at time \(t\). \(R_{cf}\) is
counterfactual reliability: it controls whether and how strongly
counterfactual evidence should influence learning.

\begin{figure}[t]
  \centering
  \resizebox{\textwidth}{!}{\begin{tikzpicture}[
  font=\scriptsize,
  module/.style={draw, rounded corners=2pt, align=center, minimum width=2.0cm, minimum height=0.7cm, fill=blue!6},
  core/.style={draw, rounded corners=2pt, align=center, minimum width=2.0cm, minimum height=0.7cm, fill=green!8},
  output/.style={draw, rounded corners=2pt, align=center, minimum width=2.0cm, minimum height=0.7cm, fill=orange!10},
  branch/.style={draw, rounded corners=2pt, align=center, minimum width=2.0cm, minimum height=0.65cm, fill=purple!7},
  arrow/.style={-{Stealth[length=2mm]}, thick},
  soft/.style={-{Stealth[length=2mm]}, thick, dashed, draw=gray!70}
]
\node[module] (Input) at (0,0) {Input interactions\\and reviews};
\node[module] (Evidence) at (2.6,0) {Evidence\\extraction};
\node[core] (EASD) at (5.2,0) {EASD\\content/style\\decomposition};
\node[core] (HSS) at (7.8,0) {HSS\\hierarchical\\style state};
\node[core] (RCR) at (10.4,0) {RCR/UCI\\reliability-gated\\routing};
\node[core] (CCV) at (13.0,0) {CCV\\controlled\\verbalizer};
\node[output] (Exp) at (15.6,0) {Target-domain\\explanation};

\node[branch] (Rating) at (7.8,-1.75) {Content-grounded\\rating branch};
\node[branch] (FCA) at (13.0,-1.75) {FCA\\faithfulness\\consistency};
\node[output] (RateOut) at (10.4,-3.0) {Target-domain\\rating};

\draw[arrow] (Input) -- (Evidence);
\draw[arrow] (Evidence) -- (EASD);
\draw[arrow] (EASD) -- (HSS);
\draw[arrow] (HSS) -- (RCR);
\draw[arrow] (RCR) -- (CCV);
\draw[arrow] (CCV) -- (Exp);
\draw[arrow] (EASD) -- node[left, font=\tiny] {$C$} (Rating);
\draw[arrow] (Rating) -- (RateOut);
\draw[arrow] (EASD) |- node[pos=0.22, left, font=\tiny] {$C$} (FCA);
\draw[arrow] (CCV) -- (FCA);
\draw[soft] (FCA) -- node[below, font=\tiny] {content alignment} (Rating);
\draw[soft] (RCR) |- node[pos=0.2, right, font=\tiny] {$R_{cf}$ controls trusted signals} (FCA);
\node[font=\tiny, text=gray!75] at (5.2,0.82) {what to explain};
\node[font=\tiny, text=gray!75] at (7.8,0.82) {how to express};
\node[font=\tiny, text=gray!75] at (10.4,0.82) {whether to trust CF};
\end{tikzpicture}
}
  \caption{ODCR overall architecture. The method maps input interactions and reviews into evidence extraction, EASD content/style decomposition, HSS hierarchical style state modeling, RCR/UCI reliability-gated counterfactual routing, CCV controlled explanation generation, and FCA faithfulness consistency. The content-grounded rating branch is separated from explanation generation while sharing content evidence.}
  \label{fig:overall_architecture}
\end{figure}

Figure~\ref{fig:overall_architecture} shows the resulting architecture. ODCR
does not ask one representation to absorb every textual factor. Instead, it
uses \(C\), \(S^{(t)}\), and \(R_{cf}\) as method-level controls for content,
expression, and counterfactual trust.

\subsection{Evidence-Anchored Content/Style Decomposition}

Evidence-Anchored Content/Style Decomposition (EASD) is the first ODCR module.
It separates stable content evidence from expression style while keeping both
anchored to observable user--item evidence. Content evidence \(C\) is tied to
rating-relevant aspects and interaction evidence. Style evidence captures
domain phrasing, emphasis, and lexical conventions that affect how content is
verbalized.

The key distinction is that ODCR does not leave shared and specific
representations to infer this split without guidance. The decomposition is
anchored: content is constrained by preference evidence, while style is
constrained by domain-specific expression signals. This turns the coarse
attribute \(W^{(t)}\) into a controllable pair: \(C\) determines what should be
explained, and \(S^{(t)}\) determines how that evidence is expressed. Rating
prediction and explanation generation can therefore share content evidence
without collapsing explanation style into the rating branch.

\subsection{Hierarchical Style State}

Hierarchical Style State (HSS) models style as more than a domain identifier.
ODCR writes the style state as
\[
  S^{(t)} = \bigl(S_{\mathrm{domain}}^{(t)}, \Delta S_{\mathrm{local}}^{(t)}\bigr),
\]
where \(S_{\mathrm{domain}}^{(t)}\) captures domain-global textual conventions
and \(\Delta S_{\mathrm{local}}^{(t)}\) captures instance-local residual style.
The domain component models broad conventions such as review tone and
domain-common phrasing. The local residual captures user--item-specific
stylistic variation that should not be forced into a single domain template.

HSS makes the D4C feedback channel more precise. If explanations feed back
into the future textual environment, ODCR localizes that feedback mainly to
\(S^{(t+1)}\). Content evidence \(C\) should remain a stable anchor for what is
being recommended and explained. This prevents dynamic feedback from
corrupting content while still allowing expression style to evolve with the
target domain.

\subsection{Reliability-Gated Counterfactual Control}

Reliability-Gated Counterfactual Control combines RCR and uncertainty-aware
counterfactual invariance (UCI). D4C shows that counterfactual domain evidence
can reduce shortcut dependence, but ODCR further recognizes that
counterfactual examples are not equally trustworthy. Some preserve content and
shift only style; others alter the recommendation evidence or introduce noisy
signals.

ODCR uses two control layers. The first layer is eligibility filtering: a
counterfactual example is usable only if it preserves the content needed for
the target explanation and does not violate the intended transfer setting. The
second layer is reliability weighting and sampling priority: usable
counterfactual examples receive different influence according to \(R_{cf}\)
and uncertainty. A compact module-level objective is
\[
  \mathcal{L}_{\mathrm{ODCR}}
  =
  \mathcal{L}_{\mathrm{fact}}
  +
  \lambda \sum_j R_{cf,j}\,\mathcal{L}_{\mathrm{cf},j},
\]
where \(\mathcal{L}_{\mathrm{fact}}\) is factual learning, \(\mathcal{L}_{\mathrm{cf},j}\)
is the counterfactual learning signal for candidate \(j\), and \(R_{cf,j}\)
controls its influence. UCI reduces trust in unstable or uncertain
counterfactual signals rather than treating every counterfactual example as
equally informative.

\begin{figure}[t]
  \centering
  \resizebox{\textwidth}{!}{\begin{tikzpicture}[
  font=\scriptsize,
  box/.style={draw, rounded corners=2pt, align=center, minimum width=2.0cm, minimum height=0.65cm, fill=white},
  gate/.style={draw, rounded corners=2pt, align=center, minimum width=2.15cm, minimum height=0.7cm, fill=blue!7},
  high/.style={draw, rounded corners=2pt, align=center, minimum width=1.75cm, minimum height=0.6cm, fill=green!10},
  med/.style={draw, rounded corners=2pt, align=center, minimum width=1.75cm, minimum height=0.6cm, fill=yellow!18},
  low/.style={draw, rounded corners=2pt, align=center, minimum width=1.75cm, minimum height=0.6cm, fill=red!8},
  arrow/.style={-{Stealth[length=2mm]}, thick},
  dline/.style={-{Stealth[length=2mm]}, thick, dashed, draw=gray!70}
]
\node[box] (CF) at (0,0) {candidate\\counterfactuals};
\node[gate] (Elig) at (2.8,0) {Layer 1\\eligibility\\filtering};
\node[box] (Use) at (5.5,0.7) {usable CF\\evidence};
\node[box] (Drop) at (5.5,-0.7) {discarded or\\held out};
\node[gate] (Score) at (8.2,0.7) {Layer 2\\reliability and\\uncertainty scoring};
\node[high] (High) at (10.8,1.45) {high\\priority};
\node[med] (Med) at (10.8,0.5) {medium\\priority};
\node[low] (Low) at (10.8,-0.45) {low\\priority};
\node[box] (Train) at (13.3,0.95) {weighted CF\\training signal};
\node[box] (Filter) at (13.3,-0.45) {down-weighted\\or filtered};
\node[box] (Gen) at (15.8,0.95) {controlled\\generator};

\draw[arrow] (CF) -- (Elig);
\draw[arrow] (Elig) -- node[above, font=\tiny] {eligible} (Use);
\draw[dline] (Elig) -- node[below, font=\tiny] {not eligible} (Drop);
\draw[arrow] (Use) -- (Score);
\draw[arrow] (Score) -- (High);
\draw[arrow] (Score) -- (Med);
\draw[arrow] (Score) -- (Low);
\draw[arrow] (High) -- (Train);
\draw[arrow] (Med) -- (Train);
\draw[dline] (Low) -- (Filter);
\draw[arrow] (Train) -- (Gen);
\draw[dline] (Filter) -- (Gen);
\node[font=\tiny, text=gray!75] at (8.2,1.65) {UCI penalizes unstable signals};
\node[font=\tiny, text=gray!75] at (13.3,1.65) {$R_{cf}$ sets sample influence};
\end{tikzpicture}
}
  \caption{RCR/UCI routing mechanism. Counterfactual candidates first pass eligibility filtering, then receive reliability and uncertainty-aware weights that determine sampling priority and training influence. Low-reliability evidence is down-weighted or filtered rather than treated as equal counterfactual supervision.}
  \label{fig:rcr_routing}
\end{figure}

Figure~\ref{fig:rcr_routing} illustrates this two-layer routing mechanism. The
figure is intended as a method diagram and as an insertion point for later
module ablations; it does not claim that a complete ablation suite is already
integrated into the current result tables.

\subsection{Controlled Causal Verbalizer and Faithfulness Alignment}

The Controlled Causal Verbalizer (CCV) is the output-side module. Instead of
learning the shortcut \(P(E\mid W)\), ODCR conditions explanation generation on
the refined variables:
\[
  P(E \mid U, I, C, S, R_{cf}) .
\]
Content evidence \(C\) controls what should be said, style state \(S\) controls
how it should be expressed, and \(R_{cf}\) controls which counterfactual
training signals are trusted. This design makes explanation generation a
controlled causal verbalization problem rather than an unconstrained
domain-language imitation problem.

Faithfulness Consistency Alignment (FCA) keeps the generated explanation tied
to content-grounded recommendation evidence. Style transfer should not
substitute for preference evidence, and a fluent target-domain explanation
should not drift away from the aspects that support the rating branch. FCA is
therefore the consistency principle linking CCV back to \(C\): the output may
adapt its expression style, but its content should remain aligned with
rating-relevant evidence.

\subsection{Implementation Mapping}

The implementation maps these method modules to the current ODCR workflow
without making workflow stages the paper's primary contribution. Preprocessing
and Step3 support EASD/HSS by preparing evidence and learning content/style
representations. Step4 supports RCR/UCI through reliability estimation,
routing, confidence buckets, uncertainty signals, and sample-weight hints.
Step5\_e supports CCV/FCA as explanation-only controlled generation. Rating
evidence remains content-grounded and separate from explanation generation,
while explanation learning consumes reliability-gated evidence for target-domain
verbalization.
